\documentclass[12pt,oneside]{book}

% ============================================================
% METADATOS
% ============================================================
\newcommand{\universidad}{Universidad de Buenos Aires (UBA)}
\newcommand{\facultad}{Facultad de Derecho}
\newcommand{\carrera}{Derecho Civil}
\newcommand{\materia}{Derechos Reales}
\newcommand{\titulo}{Apuntes de Derechos Reales}
\newcommand{\subtitulo}{Introducci\'on a la Teor\'ia General}
\newcommand{\unidadclase}{Unidad 1 --- Clase 1}
\newcommand{\catedra}{C\'atedra Prof. Farina}
\newcommand{\autor}{Isaac Edgar Camacho Ocampo}
\newcommand{\anio}{2026}

% ============================================================
% PAQUETES
% ============================================================
\usepackage[utf8]{inputenc}
\usepackage[T1]{fontenc}
\usepackage[spanish,es-nodecimaldot]{babel}

\usepackage[
    top=2.5cm,
    bottom=2.5cm,
    left=3cm,
    right=2.5cm,
    headheight=15pt
]{geometry}

\usepackage{amsmath}
\usepackage{amssymb}
\usepackage{mathtools}

\usepackage{graphicx}
\usepackage{float}
\usepackage{caption}
\usepackage{subcaption}

\usepackage{booktabs}
\usepackage{array}
\usepackage{longtable}
\usepackage{tabularx}

\usepackage{enumitem}

\usepackage{xcolor}
\usepackage[most]{tcolorbox}

\usepackage{fancyhdr}
\usepackage{ragged2e}

\usepackage[
    colorlinks=true,
    linkcolor=blue!50!black,
    citecolor=green!40!black,
    urlcolor=blue!60!black,
    pdftitle={\titulo},
    pdfauthor={\autor},
    pdfsubject={Apuntes universitarios},
    bookmarks=true
]{hyperref}

% ============================================================
% CONFIGURACIÓN
% ============================================================
\graphicspath{
    {./img/}
    {./graficos/}
    {./figuras/}
}

\setcounter{tocdepth}{2}

\setlength{\parindent}{0pt}
\setlength{\parskip}{0.5em}

\setlist[itemize]{
    topsep=0.3em,
    itemsep=0.2em
}

\setlist[enumerate]{
    topsep=0.3em,
    itemsep=0.2em
}

\clubpenalty=10000
\widowpenalty=10000

% ============================================================
% COMANDOS
% ============================================================
\newcommand{\abs}[1]{\left|#1\right|}
\newcommand{\norm}[1]{\left\lVert#1\right\rVert}
\newcommand{\vect}[1]{\mathbf{#1}}
\newcommand{\deriv}[2]{\frac{d#1}{d#2}}
\newcommand{\pderiv}[2]{\frac{\partial#1}{\partial#2}}

% ============================================================
% ENTORNOS / CAJAS
% ============================================================
\newtcolorbox{definition}[1]{
    enhanced,
    breakable,
    title={Definici\'on: #1},
    fonttitle=\bfseries,
    colback=blue!3,
    colframe=blue!50!black,
    boxrule=0.8pt,
    arc=2mm
}

\newtcolorbox{important}{
    enhanced,
    breakable,
    title={Importante},
    fonttitle=\bfseries,
    colback=yellow!8,
    colframe=orange!70!black,
    boxrule=0.8pt,
    arc=2mm
}

\newtcolorbox{example}[1]{
    enhanced,
    breakable,
    title={Ejemplo: #1},
    fonttitle=\bfseries,
    colback=green!4,
    colframe=green!40!black,
    boxrule=0.8pt,
    arc=2mm
}

\newtcolorbox{formula}[1]{
    enhanced,
    breakable,
    title={F\'ormula / Regla: #1},
    fonttitle=\bfseries,
    colback=purple!4,
    colframe=purple!50!black,
    boxrule=0.8pt,
    arc=2mm
}

\newtcolorbox{exercise}[1]{
    enhanced,
    breakable,
    title={Ejercicio: #1},
    fonttitle=\bfseries,
    colback=gray!5,
    colframe=gray!60!black,
    boxrule=0.8pt,
    arc=2mm
}

\newtcolorbox{note}{
    enhanced,
    breakable,
    title={Nota},
    fonttitle=\bfseries,
    colback=white,
    colframe=gray!50,
    boxrule=0.6pt,
    arc=2mm
}

% ============================================================
% CONFIGURACIÓN FINAL
% ============================================================
\pagestyle{fancy}
\fancyhf{}
\fancyhead[L]{\nouppercase{\leftmark}}
\fancyhead[R]{\materia}
\fancyfoot[C]{\thepage}
\renewcommand{\headrulewidth}{0.4pt}
\renewcommand{\footrulewidth}{0pt}

\fancypagestyle{plain}{
    \fancyhf{}
    \fancyfoot[C]{\thepage}
    \renewcommand{\headrulewidth}{0pt}
}

\renewcommand{\arraystretch}{1.25}

% ============================================================
% DOCUMENTO
% ============================================================
\begin{document}

% ------------------------------------------------------------
% PORTADA
% ------------------------------------------------------------
\begin{titlepage}

\thispagestyle{empty}

\begin{center}

\vspace*{1cm}

{\large \textsc{\universidad}}

\vspace{0.2cm}

{\large \textsc{\facultad}}

\vspace{0.2cm}

{\large \carrera}

\vspace{3cm}

{\Huge \textbf{\materia}}

\vspace{1cm}

{\LARGE \titulo}

\vspace{0.3cm}

{\large \subtitulo}

\vspace{0.2cm}

{\normalsize \unidadclase\ \textperiodcentered\ \catedra}

\vspace{0.8cm}

\textit{Comprender los conceptos es el primer paso para convertir el estudio en conocimiento.}

\vspace{1.2cm}

\rule{0.70\textwidth}{0.5pt}

\vspace{0.5cm}

{\large Apuntes de estudio}

\vspace{0.5cm}

\rule{0.70\textwidth}{0.5pt}

\vfill

{\large \autor}

\vspace{0.3cm}

{\large \anio}

\end{center}

\vfill

\begin{flushleft}
\textit{Este es un modesto aporte a los estudiantes de ``Derechos Reales'' no pretende sustituir el material oficial de la materia.}
\end{flushleft}

\end{titlepage}

% ------------------------------------------------------------
% ÍNDICE
% ------------------------------------------------------------
\tableofcontents

% ============================================================
% CONTENIDO
% ============================================================

% ------------------------------------------------------------
\chapter[Derechos reales y personales]{Distinci\'on entre Derechos Reales y Personales}
% ------------------------------------------------------------

\section{Concepto de derecho personal (obligacional)}

El \textbf{C\'odigo Civil y Comercial vigente} (CCyC) adopta, siguiendo la teor\'ia aplicable en la mayor\'ia de los sistemas occidentales de derecho, la distinci\'on entre \textbf{derechos reales} y \textbf{derechos personales}. Esta distinci\'on se denomina \textbf{teor\'ia cl\'asica o del v\'inculo}.

El derecho personal, tambi\'en llamado \textbf{derecho obligacional}, implica un v\'inculo jur\'idico entre sujetos. El \textbf{art\'iculo 724 del CCyC} define la obligaci\'on:

\begin{definition}{Obligaci\'on (art. 724 CCyC)}
``La obligaci\'on es una relaci\'on jur\'idica en virtud de la cual el acreedor tiene el derecho de exigir del deudor una prestaci\'on destinada a satisfacer un inter\'es l\'icito y ante el incumplimiento obtener forzadamente la satisfacci\'on de dicho inter\'es.''
\end{definition}

\subsection{Elementos del derecho personal}

Los derechos personales tienen tres elementos esenciales:

\begin{itemize}
    \item \textbf{Sujeto activo o acreedor:} quien tiene el derecho de exigir la prestaci\'on.
    \item \textbf{Sujeto pasivo o deudor:} quien est\'a obligado a cumplir la prestaci\'on.
    \item \textbf{La prestaci\'on:} conducta debida, que puede ser de dar, de hacer o de no hacer (omisi\'on).
\end{itemize}

Este v\'inculo se representa como un \textbf{v\'inculo circular} (una <<ola>>) que encierra dentro de s\'i al acreedor, al deudor y a la prestaci\'on.

\subsection{Caracteres del derecho personal}

\begin{itemize}
    \item \textbf{Relativos:} el v\'inculo existe solo entre las partes. El acreedor solo puede exigir de su deudor.
    \item \textbf{Objeto:} conducta humana (dar, hacer, no hacer).
    \item \textbf{Nacen de:} hechos o actos jur\'idicos (principalmente el contrato y la responsabilidad civil).
    \item \textbf{Prevalece} la autonom\'ia de la voluntad; la ley es supletoria.
    \item \textbf{Se extinguen} por cumplimiento o por prescripci\'on liberatoria. Son siempre temporales.
    \item No se adquieren por el mero transcurso del tiempo.
\end{itemize}

\begin{example}{obligaci\'on alimentaria y cr\'edito bancario}
En la \textbf{obligaci\'on alimentaria}, el hijo (acreedor) solo puede exigir el cumplimiento de la prestaci\'on alimentaria a quienes la ley identifica como obligados (deudores: padres, abuelos, etc.). En el \textbf{cr\'edito bancario}, el banco acreedor solo puede reclamar a su deudor.
\end{example}

\section{Concepto de derecho real}

El \textbf{art\'iculo 1882 del CCyC} define el derecho real:

\begin{definition}{Derecho real (art. 1882 CCyC)}
``El derecho real es el poder jur\'idico de estructura legal que se ejerce directamente sobre su objeto en forma aut\'onoma y que le atribuye a su titular las facultades de persecuci\'on y preferencia y las dem\'as previstas en este c\'odigo.''
\end{definition}

El derecho real implica un \textbf{v\'inculo jur\'idico directo e inmediato} entre el titular del derecho y la cosa. Ese v\'inculo es tan intenso que permite:

\begin{itemize}
    \item \textbf{Disponer de la cosa} u \textbf{obtener alg\'un beneficio} de ella en forma directa, sin intermediarios.
    \item \textbf{Perseguir la cosa} en manos de quien se encuentre (\textit{ius persequendi}).
    \item \textbf{Ser preferido} a cualquier otro titular de derecho real o personal de fecha posterior (\textit{ius preferendi}).
\end{itemize}

\subsection{Elementos del derecho real}

El derecho real tiene \textbf{solo dos elementos}: el \textbf{titular del derecho} y la \textbf{cosa}. La relaci\'on es directa e inmediata del titular con la cosa (se la representa como una flecha directa). El sujeto pasivo es toda la comunidad, que debe respetar el derecho.

\subsection{Caracteres del derecho real}

\begin{itemize}
    \item \textbf{Absolutos:} oponibles \textit{erga omnes} (a toda la comunidad).
    \item Son creados y regulados \textbf{exclusivamente por la ley} (\textit{numerus clausus}, art. 1884).
    \item \textbf{Objeto:} la cosa (bien material). Excepcionalmente, bienes inmateriales taxativamente previstos.
    \item \textbf{Nacen de:} t\'itulo suficiente y modo (tradici\'on), ley, prescripci\'on adquisitiva.
    \item Pueden ser perpetuos o temporales.
    \item Se adquieren por prescripci\'on adquisitiva.
    \item Su normativa es de \textbf{orden p\'ublico}: los particulares no pueden apartarse de ella bajo sanci\'on de nulidad.
\end{itemize}

\section{Cuadro comparativo: derechos reales vs. personales}

\begin{table}[H]
\centering
\begin{tabularx}{\textwidth}{
    >{\raggedright\arraybackslash}p{0.19\textwidth}
    >{\raggedright\arraybackslash}X
    >{\raggedright\arraybackslash}X
}
\toprule
\textbf{Aspecto} & \textbf{Derechos personales} & \textbf{Derechos reales} \\
\midrule
V\'inculo & Entre sujetos (acreedor-deudor) & Entre titular y cosa \\
Oponibilidad & Relativo (solo entre las partes) & Absoluto (\textit{erga omnes}) \\
Sujeto pasivo & Deudor determinado & Toda la comunidad \\
Objeto & Conducta humana (dar, hacer, no hacer) & La cosa (bien material) \\
Creaci\'on & Autonom\'ia de la voluntad & Exclusivamente la ley \\
Ley & Supletoria & Orden p\'ublico \\
Prescripci\'on & Liberatoria (extingue el derecho) & Adquisitiva (crea el derecho) \\
Duraci\'on & Siempre temporales & Perpetuos o temporales \\
Nacimiento del v\'inculo & Contrato, da\~no, etc. & T\'itulo y modo (tradici\'on), ley \\
\bottomrule
\end{tabularx}
\end{table}

Ambos tipos de derechos integran el \textbf{patrimonio} de una persona. Los bienes de una persona son tanto los derechos personales como los derechos reales (y las cosas materiales).

\section{Cuadro Cornell del Cap\'itulo 1}

\begin{longtable}{
    >{\raggedright\arraybackslash}p{0.26\textwidth}
    >{\raggedright\arraybackslash}p{0.66\textwidth}
}
\toprule
\textbf{Preguntas / palabras clave} & \textbf{Notas / respuestas} \\
\midrule
\endhead

\textbf{\textquestiondown Qu\'e teor\'ia adopta el CCyC para distinguir los derechos reales de los personales?} &
Adopta la teor\'ia cl\'asica o del v\'inculo, que distingue los derechos reales (v\'inculo titular-cosa) de los derechos personales (v\'inculo sujeto-sujeto). \\[0.6em]

\textbf{\textquestiondown C\'omo define el CCyC la obligaci\'on y cu\'ales son los elementos del derecho personal?} &
Art. 724: relaci\'on jur\'idica en virtud de la cual el acreedor tiene el derecho de exigir del deudor una prestaci\'on destinada a satisfacer un inter\'es l\'icito y, ante el incumplimiento, obtener forzadamente la satisfacci\'on de dicho inter\'es. Elementos: sujeto activo o acreedor, sujeto pasivo o deudor y prestaci\'on (dar, hacer o no hacer). \\[0.6em]

\textbf{\textquestiondown Qu\'e caracteres identifican a los derechos personales?} &
Son relativos; su objeto es conducta humana; nacen de hechos o actos jur\'idicos; prevalece la autonom\'ia de la voluntad y la ley es supletoria; se extinguen por cumplimiento o prescripci\'on liberatoria; son siempre temporales. \\[0.6em]

\textbf{\textquestiondown C\'omo define el art. 1882 el derecho real y cu\'ales son sus elementos?} &
Poder jur\'idico de estructura legal que se ejerce directamente sobre su objeto en forma aut\'onoma y que atribuye al titular las facultades de persecuci\'on y preferencia. Tiene dos elementos: el titular y la cosa. \\[0.6em]

\textbf{\textquestiondown Por qu\'e los derechos reales son absolutos y los personales relativos?} &
Los reales son oponibles \textit{erga omnes}: toda la comunidad debe respetar el derecho del titular. Los personales son relativos porque el v\'inculo existe solo entre las partes y el acreedor solo puede exigir de su deudor. \\[0.6em]

\textbf{\textquestiondown Qu\'e caracteres tienen los derechos reales?} &
Son absolutos; creados y regulados solo por la ley; su objeto es la cosa; nacen de t\'itulo y modo, ley o prescripci\'on adquisitiva; pueden ser perpetuos o temporales; su normativa es de orden p\'ublico. \\[0.6em]

\textbf{\textquestiondown Cu\'al es la diferencia esencial entre un derecho real y uno personal? \textquestiondown Qu\'e integra el patrimonio?} &
El personal es un v\'inculo entre sujetos (relativo); el real es un v\'inculo directo entre titular y cosa (absoluto, \textit{erga omnes}). El patrimonio est\'a integrado tanto por derechos personales como por derechos reales. \\

\midrule
\multicolumn{2}{p{\dimexpr0.92\textwidth+2\tabcolsep\relax}}{%
\textbf{Resumen:} El CCyC adopta la teor\'ia cl\'asica o del v\'inculo. El derecho personal vincula a un acreedor y a un deudor en torno a una prestaci\'on y es relativo, temporal y regido por la autonom\'ia de la voluntad. El derecho real vincula directamente al titular con la cosa, es absoluto (\textit{erga omnes}), solo la ley lo crea y su normativa es de orden p\'ublico.} \\
\bottomrule
\end{longtable}

% ------------------------------------------------------------
\chapter[Regulaci\'on en el CCyC]{Regulaci\'on en el C\'odigo Civil y Comercial}
% ------------------------------------------------------------

\section{Estructura metodol\'ogica del CCyC}

El C\'odigo Civil y Comercial tiene una \textbf{metodolog\'ia novedosa} respecto del anterior C\'odigo de V\'elez S\'arsfield.

\subsection{T\'itulo Preliminar}

Contiene los principios rectores aplicables a todo el C\'odigo Civil y Comercial:

\begin{itemize}
    \item Principio de buena fe.
    \item Abuso del derecho.
    \item Abuso de posici\'on dominante.
    \item Preeminencia del orden p\'ublico.
    \item Tratados de derechos humanos como fuente y gu\'ia de interpretaci\'on.
    \item \textbf{Cap\'itulo 4:} Bienes y cosas (fundamental para los derechos reales).
\end{itemize}

\subsection{Libros del C\'odigo}

\begin{table}[H]
\centering
\begin{tabularx}{\textwidth}{
    >{\raggedright\arraybackslash}p{0.2\textwidth}
    >{\raggedright\arraybackslash}p{0.27\textwidth}
    >{\raggedright\arraybackslash}X
}
\toprule
\textbf{Libro} & \textbf{Denominaci\'on} & \textbf{Contenido relevante} \\
\midrule
T\'itulo Preliminar & Principios rectores & Buena fe, orden p\'ublico, derechos humanos \\
Libro Primero & Parte General & Persona, capacidad \\
Libro Segundo & Relaciones de Familia & Derecho de familia \\
Libro Tercero & Derechos Personales & Obligaciones, contratos \\
Libro Cuarto & Derechos Reales & Teor\'ia general y derechos reales en particular \\
Libro Quinto & Transmisi\'on por causa de muerte & Sucesiones \\
T\'itulo Sexto & Disposiciones Comunes & A derechos reales y personales \\
\bottomrule
\end{tabularx}
\end{table}

\begin{important}
El estudio de los derechos reales \textbf{excede el Libro Cuarto}: es necesario integrar el T\'itulo Preliminar, el T\'itulo Sexto y el resto del ordenamiento.
\end{important}

\section{Art\'iculos clave: 15, 16 y 17 del CCyC}

\begin{formula}{Art. 15 CCyC --- Titularidad}
``Las personas son titulares de los derechos individuales sobre los bienes que integran su patrimonio conforme con lo que se establece en este C\'odigo.''

El patrimonio est\'a integrado por derechos personales y derechos reales.
\end{formula}

\begin{formula}{Art. 16 CCyC --- Bienes y cosas}
``Los derechos pueden recaer sobre bienes susceptibles de valor econ\'omico. Cuando son materiales se llaman cosas. Las disposiciones sobre las cosas se aplican a la energ\'ia y a las fuerzas naturales susceptibles de ser puestas al servicio del hombre.''

Esta norma es fundamental para determinar el \textbf{objeto} del derecho real.
\end{formula}

\begin{formula}{Art. 17 CCyC --- Derechos sobre el cuerpo humano}
Los derechos sobre el cuerpo humano y sus partes tienen valor afectivo, cient\'ifico, terap\'eutico, humanitario o social, pero no valor comercial. Se rigen por leyes especiales (\'organos, genes, etc.).

Las partes del cuerpo no son objeto de derechos reales por carecer de valor econ\'omico.
\end{formula}

\section{Cuadro Cornell del Cap\'itulo 2}

\begin{longtable}{
    >{\raggedright\arraybackslash}p{0.26\textwidth}
    >{\raggedright\arraybackslash}p{0.66\textwidth}
}
\toprule
\textbf{Preguntas / palabras clave} & \textbf{Notas / respuestas} \\
\midrule
\endhead

\textbf{\textquestiondown Qu\'e contiene el T\'itulo Preliminar del CCyC?} &
Los principios rectores aplicables a todo el C\'odigo: buena fe, abuso del derecho, abuso de posici\'on dominante, preeminencia del orden p\'ublico y tratados de derechos humanos como fuente y gu\'ia de interpretaci\'on. Incluye el Cap\'itulo 4, sobre bienes y cosas. \\[0.6em]

\textbf{\textquestiondown C\'omo se estructura el CCyC?} &
T\'itulo Preliminar; Libro Primero (Parte General); Libro Segundo (Familia); Libro Tercero (Derechos Personales); Libro Cuarto (Derechos Reales); Libro Quinto (Transmisi\'on por causa de muerte); T\'itulo Sexto (Disposiciones Comunes). \\[0.6em]

\textbf{\textquestiondown En qu\'e libro est\'an los derechos reales? \textquestiondown Es suficiente estudiar solo ese libro?} &
Est\'an en el Libro Cuarto, pero su estudio lo excede: hay que integrar el T\'itulo Preliminar, el T\'itulo Sexto y el resto del ordenamiento. \\[0.6em]

\textbf{\textquestiondown Qu\'e establece el art. 15?} &
Las personas son titulares de derechos individuales sobre los bienes que integran su patrimonio. \\[0.6em]

\textbf{\textquestiondown Cu\'ando un bien es una cosa seg\'un el art. 16?} &
Los derechos pueden recaer sobre bienes susceptibles de valor econ\'omico; cuando son materiales se llaman cosas. Las disposiciones sobre las cosas se aplican tambi\'en a la energ\'ia y a las fuerzas naturales. Es fundamental para determinar el objeto del derecho real. \\[0.6em]

\textbf{\textquestiondown Por qu\'e las partes del cuerpo humano no son objeto de derechos reales (art. 17)?} &
Porque tienen valor afectivo, cient\'ifico, terap\'eutico, humanitario o social, pero no valor comercial ni econ\'omico; se rigen por leyes especiales. \\

\midrule
\multicolumn{2}{p{\dimexpr0.92\textwidth+2\tabcolsep\relax}}{%
\textbf{Resumen:} El CCyC contiene un T\'itulo Preliminar con principios rectores, libros tem\'aticos y un T\'itulo Sexto de disposiciones comunes. Los derechos reales se regulan en el Libro Cuarto, pero su estudio requiere integrar otras partes del C\'odigo. Los arts. 15, 16 y 17 establecen la titularidad sobre los bienes del patrimonio, el concepto de cosa y la exclusi\'on del cuerpo humano como objeto de derechos reales.} \\
\bottomrule
\end{longtable}

% ------------------------------------------------------------
\chapter[Teor\'ia general de los derechos reales]{Teor\'ia General de los Derechos Reales}
% ------------------------------------------------------------

\section{Principios comunes (arts. 1882 a 1891)}

Los art\'iculos \textbf{1882 a 1891 del CCyC} regulan los principios comunes a todos los derechos reales. Esta parte general reviste \textbf{importancia fundamental}:

\begin{itemize}
    \item Es la base te\'orica de toda la materia.
    \item Muchos de sus principios no se repiten al estudiar cada derecho real en particular.
\end{itemize}

\section{Numerus clausus (art. 1884 CCyC)}

El \textbf{art\'iculo 1884 CCyC} consagra el principio de \textit{numerus clausus} (n\'umeros cerrados):

\begin{formula}{Art. 1884 CCyC --- Estructura de los derechos reales}
``La regulaci\'on de los derechos reales en cuanto a sus elementos, contenido, adquisici\'on, constituci\'on, modificaci\'on, transmisi\'on, duraci\'on, extinci\'on (...) est\'a establecida solo por la ley. Es nula la configuraci\'on de un derecho real no previsto en la ley o la modificaci\'on de su estructura.''
\end{formula}

\begin{itemize}
    \item \textbf{Todo} en materia de derechos reales es regulado por la ley: elementos, contenido, adquisici\'on, constituci\'on, modificaci\'on, transmisi\'on, duraci\'on y extinci\'on.
    \item Si los particulares se apartan de lo regulado por la ley, la consecuencia es la \textbf{nulidad}.
    \item El derecho real es un se\~nor\'io de la voluntad que se ejerce en forma \textbf{aut\'onoma e independiente} de otra voluntad.
\end{itemize}

\section{Convalidaci\'on (art. 1885 CCyC)}

El \textbf{principio de convalidaci\'on} responde al principio general de \textit{nemo plus iuris}.

\begin{important}
\textbf{\textit{Nemo plus iuris}:} nadie puede transmitir a otro un derecho mejor o m\'as extenso que el que ten\'ia. El acto nulo se convalida si el transmitente adquiere el derecho posteriormente.
\end{important}

\begin{formula}{Art. 1885 CCyC --- Convalidaci\'on}
``Si alguien constituye o transmite un derecho que no tiene, y posteriormente lo adquiere, la constituci\'on o transmisi\'on queda convalidada.''
\end{formula}

\begin{example}{transmisi\'on de una cosa sin ser propietario}
Quien tiene en su poder una cosa pero \textbf{no tiene la propiedad} de ella, y la transmite, no puede hacer que quien la recibe adquiera un derecho de propiedad (no puede transmitir un derecho mejor que el que ten\'ia). Sin embargo, si el transmitente adquiere la propiedad posteriormente, ese acto queda convalidado.
\end{example}

\section{Persecuci\'on y preferencia}

Son las dos facultades caracter\'isticas del derecho real.

\begin{definition}{\textit{Ius persequendi} (derecho de persecuci\'on)}
El titular puede \textbf{perseguir la cosa} en poder de quien se encuentre. No importa si la cosa cambi\'o de titular: el derecho la sigue. Se lo representa como una flecha que persigue la cosa sin importar qui\'en la tenga.
\end{definition}

\begin{definition}{\textit{Ius preferendi} (derecho de preferencia)}
El titular es \textbf{preferido} a cualquier otro que pudiera tener sobre la cosa alg\'un derecho real o personal de fecha posterior.
\end{definition}

\section{Los 14 derechos reales (art. 1887 CCyC)}

El \textbf{art\'iculo 1887 del CCyC} enumera taxativamente los 14 derechos reales.

% TODO: contenido incompleto o ambiguo en las notas originales.
% La tabla de las notas lista 15 filas (incluye la numeraci\'on "14b" para Prenda) y repite
% "Superficie" (con y sin propiedad superficiaria); se conserva tal como figura en el apunte.
\begin{table}[H]
\centering
\begin{tabularx}{\textwidth}{
    >{\centering\arraybackslash}p{0.08\textwidth}
    >{\raggedright\arraybackslash}X
    >{\raggedright\arraybackslash}X
}
\toprule
\textbf{N.\textsuperscript{o}} & \textbf{Derecho real} & \textbf{Sobre cosa} \\
\midrule
1 & Dominio & Propia \\
2 & Condominio & Propia (varios titulares) \\
3 & Propiedad horizontal & Propia \\
4 & Conjuntos inmobiliarios & Propia \\
5 & Tiempo compartido & Propia \\
6 & Cementerio privado & Propia \\
7 & Superficie (con propiedad superficiaria) & Propia \\
8 & Usufructo & Ajena \\
9 & Uso & Ajena \\
10 & Habitaci\'on & Ajena \\
11 & Servidumbre & Ajena \\
12 & Superficie (sin propiedad superficiaria) & Ajena \\
13 & Hipoteca & Ajena (garant\'ia) \\
14 & Anticresis & Ajena (garant\'ia) \\
14b & Prenda & Ajena (garant\'ia) \\
\bottomrule
\end{tabularx}
\end{table}

\section{Clasificaci\'on de los derechos reales (art. 1888 CCyC)}

\subsection{Sobre cosa propia vs. cosa ajena}

\begin{table}[H]
\centering
\begin{tabularx}{\textwidth}{XX}
\toprule
\textbf{Sobre cosa propia} & \textbf{Sobre cosa ajena} \\
\midrule
Dominio & Usufructo \\
Condominio & Uso \\
Propiedad horizontal & Habitaci\'on \\
Conjuntos inmobiliarios & Servidumbre \\
Tiempo compartido & Superficie (sin propiedad superficiaria) \\
Cementerio privado & Hipoteca \\
Superficie (con propiedad superficiaria) & Anticresis / Prenda \\
\bottomrule
\end{tabularx}
\end{table}

Los derechos sobre cosa ajena constituyen una carga o gravamen para el titular del dominio.

\begin{example}{dominio e hipoteca}
Una persona es titular de dominio (derecho real sobre cosa propia) y constituye una hipoteca a favor del banco. El banco adquiere un derecho real sobre la cosa de esa persona (la hipoteca es un derecho real sobre cosa ajena).
\end{example}

\subsection{Principales vs. accesorios}

\begin{itemize}
    \item \textbf{Principales:} derechos sobre cosa propia y derechos de disfrute sobre cosa ajena (usufructo, servidumbre, uso, habitaci\'on, superficie). Son independientes.
    \item \textbf{Accesorios:} derechos reales de garant\'ia (hipoteca, prenda, anticresis). Dependen de un cr\'edito al que acceden y siguen al cr\'edito.
\end{itemize}

\subsection{Registrables vs. no registrables}

Seg\'un si su adquisici\'on, transmisi\'on o extinci\'on requieren inscripci\'on en un registro p\'ublico:

\begin{itemize}
    \item \textbf{Registrables:} inmuebles (Registro de la Propiedad Inmueble) y automotores.
    \item \textbf{No registrables:} cosas muebles en general.
\end{itemize}

\subsection{Ejercidos por posesi\'on vs. cuasiposesi\'on}

Algunos derechos reales se ejercen mediante la \textbf{posesi\'on} (poder f\'isico sobre la cosa) y otros mediante la \textbf{cuasiposesi\'on}. Este tema se desarrolla en profundidad en la clase siguiente.

\section{Cuadro Cornell del Cap\'itulo 3}

\begin{longtable}{
    >{\raggedright\arraybackslash}p{0.26\textwidth}
    >{\raggedright\arraybackslash}p{0.66\textwidth}
}
\toprule
\textbf{Preguntas / palabras clave} & \textbf{Notas / respuestas} \\
\midrule
\endhead

\textbf{\textquestiondown Qu\'e son los principios comunes de los derechos reales (arts. 1882-1891)?} &
Reglas generales aplicables a todos los derechos reales: definici\'on, estructura legal, convalidaci\'on, persecuci\'on, preferencia y enumeraci\'on taxativa. Son la base te\'orica de la materia y no se repiten al estudiar cada derecho real. \\[0.6em]

\textbf{\textquestiondown Qu\'e significa \textit{numerus clausus} y cu\'al es su consecuencia (art. 1884)?} &
Solo la ley regula los derechos reales (elementos, contenido, adquisici\'on, constituci\'on, modificaci\'on, transmisi\'on, duraci\'on y extinci\'on). Los particulares no pueden crear derechos reales nuevos ni modificar su estructura; si lo hacen, la sanci\'on es la nulidad. \\[0.6em]

\textbf{\textquestiondown Qu\'e dice el \textit{nemo plus iuris}?} &
Nadie puede transmitir a otro un derecho mejor o m\'as extenso que el que tiene. Es la base del principio de convalidaci\'on. \\[0.6em]

\textbf{\textquestiondown Qu\'e es la convalidaci\'on (art. 1885)? Ejemplo.} &
Si alguien constituye o transmite un derecho que no tiene y luego lo adquiere, la constituci\'on o transmisi\'on queda convalidada. Ejemplo: quien tiene una cosa sin ser propietario y la transmite no hace adquirir propiedad al receptor; si luego adquiere la propiedad, el acto queda convalidado. \\[0.6em]

\textbf{\textquestiondown Qu\'e son el \textit{ius persequendi} y el \textit{ius preferendi}?} &
Facultades caracter\'isticas del derecho real. Persecuci\'on: perseguir la cosa en poder de quien se encuentre. Preferencia: ser preferido a cualquier otro titular de derecho real o personal de fecha posterior. \\[0.6em]

\textbf{\textquestiondown Cu\'ales son los 14 derechos reales del art. 1887?} &
Dominio, condominio, propiedad horizontal, conjuntos inmobiliarios, tiempo compartido, cementerio privado, superficie, usufructo, uso, habitaci\'on, servidumbre, hipoteca, anticresis y prenda. \\[0.6em]

\textbf{\textquestiondown C\'omo distingue el art. 1888 entre cosa propia y ajena?} &
Propia: dominio, condominio, propiedad horizontal, etc. Ajena: usufructo, servidumbre, hipoteca, etc., que constituyen carga o gravamen para el titular del dominio. \\[0.6em]

\textbf{\textquestiondown Qu\'e diferencia hay entre derechos reales principales y accesorios?} &
Los principales (sobre cosa propia y de disfrute sobre cosa ajena) son independientes. Los accesorios son los de garant\'ia (hipoteca, prenda, anticresis), dependen de un cr\'edito y siguen al cr\'edito. \\[0.6em]

\textbf{\textquestiondown Qu\'e son los derechos registrables y no registrables?} &
Registrables: requieren inscripci\'on en un registro p\'ublico (inmuebles, automotores). No registrables: cosas muebles en general. \\[0.6em]

\textbf{\textquestiondown Qu\'e distingue la posesi\'on de la cuasiposesi\'on?} &
Algunos derechos reales se ejercen por posesi\'on (poder f\'isico sobre la cosa) y otros por cuasiposesi\'on; el tema se desarrolla en la clase siguiente. \\

\midrule
\multicolumn{2}{p{\dimexpr0.92\textwidth+2\tabcolsep\relax}}{%
\textbf{Resumen:} Los arts. 1882 a 1891 contienen los principios comunes a todos los derechos reales. Solo la ley crea y regula los derechos reales (\textit{numerus clausus}, art. 1884), con sanci\'on de nulidad. El art. 1885 convalida la transmisi\'on de un derecho que luego se adquiere, en armon\'ia con el \textit{nemo plus iuris}. Las facultades de persecuci\'on y preferencia son propias del derecho real. El art. 1887 enumera taxativamente los derechos reales, que el art. 1888 clasifica en sobre cosa propia o ajena, principales o accesorios, registrables o no registrables, y ejercidos por posesi\'on o cuasiposesi\'on.} \\
\bottomrule
\end{longtable}

\section{Cuadro Cornell general de la clase}

\begin{longtable}{
    >{\raggedright\arraybackslash}p{0.26\textwidth}
    >{\raggedright\arraybackslash}p{0.66\textwidth}
}
\toprule
\textbf{Preguntas / palabras clave} & \textbf{Notas / respuestas} \\
\midrule
\endhead

\textbf{\textquestiondown Cu\'al es la diferencia clave entre derecho personal y derecho real?} &
Personal: v\'inculo entre sujetos (relativo). Real: v\'inculo directo titular-cosa (absoluto, \textit{erga omnes}). \\[0.6em]

\textbf{Art. 724 y art. 1882: definiciones} &
Obligaci\'on: relaci\'on jur\'idica por la cual el acreedor puede exigir al deudor una prestaci\'on para satisfacer un inter\'es l\'icito (tres elementos: acreedor, deudor y prestaci\'on). Derecho real: poder jur\'idico de estructura legal, ejercido directamente sobre la cosa, aut\'onomo, que confiere persecuci\'on y preferencia. \\[0.6em]

\textbf{\textquestiondown Qu\'e establecen los arts. 15, 16 y 17?} &
Art. 15: titularidad de derechos sobre los bienes del patrimonio. Art. 16: los bienes materiales se llaman cosas; se aplica tambi\'en a la energ\'ia y a las fuerzas naturales. Art. 17: el cuerpo humano y sus partes tienen valor afectivo, no comercial, y no son objeto de derechos reales. \\[0.6em]

\textbf{\textit{Numerus clausus} y convalidaci\'on} &
Solo la ley crea derechos reales (sanci\'on: nulidad). Quien transmite lo que no tiene y luego lo adquiere convalida el acto; el fundamento es el \textit{nemo plus iuris}. \\[0.6em]

\textbf{\textquestiondown C\'omo se clasifican los derechos reales?} &
Sobre cosa propia o ajena; principales o accesorios; registrables o no registrables; ejercidos por posesi\'on o cuasiposesi\'on. \\

\midrule
\multicolumn{2}{p{\dimexpr0.92\textwidth+2\tabcolsep\relax}}{%
\textbf{Resumen:} Los derechos reales son poderes jur\'idicos de estructura legal que vinculan directamente a una persona con una cosa, son absolutos (\textit{erga omnes}) y solo pueden ser creados por la ley (\textit{numerus clausus}). Se distinguen de los derechos personales, que vinculan a dos sujetos (acreedor-deudor) y son relativos. El CCyC regula 14 derechos reales, clasificados en principales y accesorios, sobre cosa propia o ajena, registrables o no. Los principios de persecuci\'on, preferencia, convalidaci\'on y \textit{nemo plus iuris} son aplicables a todos.} \\
\bottomrule
\end{longtable}

% ------------------------------------------------------------
\chapter[Ejercicios]{Ejercicios de evaluaci\'on}
% ------------------------------------------------------------

Las siguientes preguntas abarcan todo el contenido de la clase y est\'an organizadas por niveles: \textbf{b\'asico} (fundamentos), \textbf{intermedio} (aplicaci\'on) y \textbf{avanzado} (an\'alisis y casos).

\section{Nivel b\'asico}

\begin{exercise}{N.\textsuperscript{o} 1}
\textbf{Consigna.} Defina el concepto de obligaci\'on seg\'un el art\'iculo 724 del CCyC e identifique sus tres elementos.

\textbf{Respuesta.} La obligaci\'on es una relaci\'on jur\'idica en virtud de la cual el acreedor tiene el derecho de exigir del deudor una prestaci\'on destinada a satisfacer un inter\'es l\'icito y, ante el incumplimiento, obtener forzadamente la satisfacci\'on de dicho inter\'es. Sus tres elementos son: (1) sujeto activo o acreedor, (2) sujeto pasivo o deudor, y (3) la prestaci\'on (dar, hacer o no hacer).
\end{exercise}

\begin{exercise}{N.\textsuperscript{o} 2}
\textbf{Consigna.} \textquestiondown Cu\'al es la diferencia esencial entre un derecho personal y un derecho real? Explique con un ejemplo.

\textbf{Respuesta.} La diferencia esencial es la naturaleza del v\'inculo. En el derecho personal el v\'inculo es entre dos sujetos (acreedor y deudor) y es relativo. En el derecho real el v\'inculo es directo e inmediato entre el titular y la cosa, y es absoluto (\textit{erga omnes}). Ejemplo: si se presta dinero a Juan (derecho personal), solo Juan debe; si se es due\~no de un inmueble (derecho real), toda la comunidad debe respetar esa propiedad.
\end{exercise}

\begin{exercise}{N.\textsuperscript{o} 3}
\textbf{Consigna.} \textquestiondown Por qu\'e los derechos personales son ``relativos'' y los derechos reales son ``absolutos''?

\textbf{Respuesta.} Los derechos personales son relativos porque el v\'inculo existe solo entre las partes: el acreedor solo puede exigir el cumplimiento de su propio deudor. Los derechos reales son absolutos porque son oponibles \textit{erga omnes}: toda la comunidad (no solo una persona determinada) tiene la obligaci\'on de respetar el derecho del titular.
\end{exercise}

\begin{exercise}{N.\textsuperscript{o} 4}
\textbf{Consigna.} Defina el derecho real seg\'un el art\'iculo 1882 del CCyC e identifique sus elementos.

\textbf{Respuesta.} Seg\'un el art. 1882 CCyC, el derecho real es el poder jur\'idico de estructura legal que se ejerce directamente sobre su objeto en forma aut\'onoma y que le atribuye a su titular las facultades de persecuci\'on y preferencia y las dem\'as previstas en el C\'odigo. Sus elementos son: (1) el titular del derecho y (2) la cosa (objeto).
\end{exercise}

\begin{exercise}{N.\textsuperscript{o} 5}
\textbf{Consigna.} \textquestiondown Qu\'e integra el patrimonio de una persona seg\'un el CCyC?

\textbf{Respuesta.} El patrimonio de una persona est\'a integrado por bienes, que incluyen tanto los derechos personales como los derechos reales. Ambos tipos de derechos son bienes patrimoniales.
\end{exercise}

\begin{exercise}{N.\textsuperscript{o} 6}
\textbf{Consigna.} \textquestiondown Qu\'e establece el art\'iculo 16 del CCyC? \textquestiondown Por qu\'e es fundamental para los derechos reales?

\textbf{Respuesta.} El art\'iculo 16 CCyC establece que los derechos pueden recaer sobre bienes susceptibles de valor econ\'omico y que, cuando son materiales, se llaman cosas; tambi\'en se aplica a la energ\'ia y a las fuerzas naturales. Es fundamental porque define el objeto de los derechos reales: la cosa (bien material con valor econ\'omico).
\end{exercise}

\section{Nivel intermedio}

\begin{exercise}{N.\textsuperscript{o} 7}
\textbf{Consigna.} \textquestiondown Qu\'e significa el principio de \textit{numerus clausus} en materia de derechos reales? \textquestiondown Cu\'al es su consecuencia jur\'idica?

\textbf{Respuesta.} El \textit{numerus clausus} (n\'umeros cerrados) significa que solo la ley puede crear derechos reales. Los particulares no pueden crear nuevos derechos reales ni modificar el contenido de los existentes. La consecuencia jur\'idica es la nulidad: si los particulares configuran un derecho real no previsto en la ley o modifican su estructura, el acto es nulo (art. 1884 CCyC).
\end{exercise}

\begin{exercise}{N.\textsuperscript{o} 8}
\textbf{Consigna.} Explique el principio de convalidaci\'on (art. 1885 CCyC) y su relaci\'on con el \textit{nemo plus iuris}.

\textbf{Respuesta.} El \textit{nemo plus iuris} establece que nadie puede transmitir a otro un derecho mejor o m\'as extenso que el que ten\'ia. Entonces, si alguien transmite un derecho que no tiene, el acto nace nulo. Sin embargo, el principio de convalidaci\'on (art. 1885 CCyC) establece que si el transmitente adquiere ese derecho con posterioridad a la transmisi\'on, el acto queda autom\'aticamente convalidado.
\end{exercise}

\begin{exercise}{N.\textsuperscript{o} 9}
\textbf{Consigna.} \textquestiondown Qu\'e es el \textit{ius persequendi} y el \textit{ius preferendi}? \textquestiondown Por qu\'e son facultades caracter\'isticas del derecho real?

\textbf{Respuesta.} El \textit{ius persequendi} es la facultad del titular de perseguir la cosa en poder de quien se encuentre, sin importar qui\'en la tenga actualmente. El \textit{ius preferendi} es la facultad de ser preferido frente a otros titulares de derechos reales o personales de fecha posterior. Son caracter\'isticas del derecho real porque derivan de la relaci\'on directa e inmediata entre el titular y la cosa, y de su car\'acter absoluto (\textit{erga omnes}).
\end{exercise}

\begin{exercise}{N.\textsuperscript{o} 10}
\textbf{Consigna.} Enumere los 14 derechos reales previstos en el art\'iculo 1887 del CCyC.

\textbf{Respuesta.} Los 14 derechos reales son: (1) dominio, (2) condominio, (3) propiedad horizontal, (4) conjuntos inmobiliarios, (5) tiempo compartido, (6) cementerio privado, (7) superficie, (8) usufructo, (9) uso, (10) habitaci\'on, (11) servidumbre, (12) hipoteca, (13) anticresis y (14) prenda.
\end{exercise}

\begin{exercise}{N.\textsuperscript{o} 11}
\textbf{Consigna.} \textquestiondown Cu\'al es la diferencia entre derechos reales sobre cosa propia y sobre cosa ajena? D\'e un ejemplo de cada uno.

\textbf{Respuesta.} Los derechos reales sobre cosa propia son aquellos en que el titular ejerce el derecho sobre su propia cosa (por ejemplo, el dominio: ser due\~no de una casa). Los derechos reales sobre cosa ajena son aquellos en que el titular ejerce el derecho sobre la cosa de otro, constituyendo una carga o gravamen para el due\~no (por ejemplo, la hipoteca: el banco tiene un derecho real sobre el inmueble en garant\'ia de un cr\'edito).
\end{exercise}

\begin{exercise}{N.\textsuperscript{o} 12}
\textbf{Consigna.} Diferencie los derechos reales principales de los accesorios, con ejemplos.

\textbf{Respuesta.} Los derechos reales principales son independientes: existen por s\'i mismos (por ejemplo, dominio, usufructo, servidumbre). Los derechos reales accesorios son los de garant\'ia: existen en funci\'on de un cr\'edito al que acceden (por ejemplo, hipoteca, prenda, anticresis). Si se extingue el cr\'edito principal, se extingue tambi\'en el derecho real de garant\'ia.
\end{exercise}

\section{Nivel avanzado}

\begin{exercise}{N.\textsuperscript{o} 13}
\textbf{Consigna.} \textquestiondown Por qu\'e la normativa en materia de derechos reales es de orden p\'ublico? \textquestiondown Qu\'e consecuencias tiene esto?

\textbf{Respuesta.} Es de orden p\'ublico porque los derechos reales trascienden el inter\'es particular: su ejercicio afecta a toda la comunidad. Las consecuencias son: (a) los particulares no pueden crear nuevos derechos reales, (b) no pueden modificar el contenido de los existentes, (c) la violaci\'on de estas normas acarrea nulidad y (d) no se aplica la autonom\'ia de la voluntad como en los contratos.
\end{exercise}

\begin{exercise}{N.\textsuperscript{o} 14}
\textbf{Consigna.} \textquestiondown Qu\'e establece el art\'iculo 17 del CCyC sobre el cuerpo humano? \textquestiondown Por qu\'e las partes del cuerpo no son objeto de derechos reales?

\textbf{Respuesta.} El art\'iculo 17 CCyC establece que los derechos sobre el cuerpo humano o sus partes no tienen valor comercial, sino valor afectivo, cient\'ifico, terap\'eutico, humanitario o social, y se rigen por leyes especiales. No son objeto de derechos reales porque el art\'iculo 16 exige valor econ\'omico para que un bien sea considerado ``cosa'', y las partes del cuerpo carecen de ese valor econ\'omico.
\end{exercise}

\begin{exercise}{N.\textsuperscript{o} 15}
\textbf{Consigna.} Analice la siguiente situaci\'on: Pedro vende un inmueble a Mar\'ia en enero, pero Pedro no era el propietario en ese momento. En marzo, Pedro adquiere la propiedad del inmueble. \textquestiondown Qu\'e ocurre con la venta realizada en enero? Fundamente con la norma aplicable.

\textbf{Respuesta.} Conforme al principio de \textit{nemo plus iuris}, Pedro no pod\'ia transmitir a Mar\'ia un derecho que no ten\'ia. La venta de enero naci\'o nula en cuanto a la transmisi\'on del derecho real. Sin embargo, en virtud del principio de convalidaci\'on previsto en el art\'iculo 1885 del CCyC, al adquirir Pedro la propiedad en marzo, la transmisi\'on realizada en enero queda autom\'aticamente convalidada. Mar\'ia adquiere el dominio retroactivamente.
\end{exercise}

% ------------------------------------------------------------
\chapter[Material complementario: m\'etodo Feynman]{Material complementario: explicaciones did\'acticas (m\'etodo Feynman)}
% ------------------------------------------------------------

Las siguientes explicaciones reformulan los conceptos centrales de la clase con lenguaje cotidiano y ejemplos sencillos.

\section{Derechos personales}

Supongamos que se prestan \$10.000 a un amigo, Juan. En ese momento nace una \textbf{obligaci\'on}: quien presta es el acreedor, Juan es el deudor y la prestaci\'on es devolver el dinero.

El dinero solo puede reclamarse a Juan; no puede reclamarse a Mar\'ia, aunque sea la novia de Juan. Eso es lo que significa que el derecho personal sea \textbf{relativo}: solo une a dos personas espec\'ificas.

Si Juan no paga, el acreedor puede acudir a la justicia y obligarlo a pagar. Pero cuando cobra, el v\'inculo se termina: es una relaci\'on temporal, entre personas.

\begin{note}
\textbf{Frase clave.} Derecho personal = v\'inculo entre personas = relativo = temporal.
\end{note}

\section{Derechos reales}

Supongamos que se compra un departamento, se firma la escritura y se realiza la entrega (tradici\'on). Desde ese momento, el departamento es del comprador, que tiene un \textbf{derecho real} sobre esa cosa.

Esto significa que \textbf{todo el mundo} tiene que respetar que ese departamento le pertenece. No solo el vendedor: cualquier persona que quiera ocuparlo o reclamarlo tiene que enfrentarse al derecho del titular. Eso es lo que significa \textbf{absoluto} u \textit{erga omnes}.

Si el departamento fuera robado, el titular puede \textbf{perseguirlo}: no importa qui\'en lo tenga ahora, el derecho de dominio ``sigue'' a la cosa. Eso es el \textit{ius persequendi}.

\begin{note}
\textbf{Frase clave.} Derecho real = v\'inculo entre persona y cosa = absoluto = \textit{erga omnes} = se persigue la cosa dondequiera que est\'e.
\end{note}

\section{Numerus clausus}

Los derechos reales son como los planetas del sistema solar: solo existen los que la ley nombra, y no puede inventarse un planeta nuevo.

Si dos personas en un contrato inventan un nuevo ``derecho real'' que no existe en la ley (por ejemplo, un ``super-dominio compartido especial''), ese acuerdo es \textbf{nulo} en cuanto al derecho real: la ley no lo reconoce. Esto es as\'i porque los derechos reales son de \textbf{orden p\'ublico}: importan a toda la sociedad, no solo a las partes del contrato.

\begin{note}
\textbf{Frase clave.} Los derechos reales son un cat\'alogo cerrado. Solo existen los 14 que la ley crea. Inventar uno = nulidad.
\end{note}

\section{Convalidaci\'on y nemo plus iuris}

Regla b\'asica: no se puede dar algo que no se tiene. Si alguien no es due\~no de un auto, no puede venderlo como si lo fuera; si lo hace, ese acto es nulo.

Pero hay una salvedad: si despu\'es esa persona se convierte en due\~na (por ejemplo, el due\~no original se lo vende), la venta anterior queda \textbf{autom\'aticamente convalidada}, como si hubiera sido v\'alida desde el principio. Esto existe para proteger a quien compr\'o de buena fe y para dar seguridad jur\'idica al sistema.

\begin{note}
\textbf{Frase clave.} \textit{Nemo plus iuris}: no se transmite m\'as derecho del que se tiene. Si se adquiere despu\'es, el acto anterior se convalida solo.
\end{note}

\section{Los 14 derechos reales}

Los 14 derechos reales pueden pensarse como 14 tipos de relaciones distintas que se pueden tener con una cosa:

\begin{itemize}
    \item \textbf{Sobre la propia cosa:} dominio (due\~no absoluto), condominio (due\~no junto con otros), propiedad horizontal (due\~no de un departamento en un edificio), etc.
    \item \textbf{Sobre la cosa de otro:} usufructo (usar y disfrutar la cosa de otro sin ser due\~no), servidumbre (derecho de paso por el fundo ajeno), hipoteca (el banco tiene derecho sobre el inmueble si la deuda no se paga), etc.
\end{itemize}

Cada uno tiene reglas propias definidas por la ley; no pueden mezclarse ni inventarse variantes.

\begin{note}
\textbf{Frase clave.} 14 tipos de relaciones jur\'idicas con las cosas. Unos son sobre la propia cosa, otros sobre la cosa de otro. Todos creados por ley.
\end{note}

\end{document}
